% ======================================================================
% DRAFT v2 — Stellar Mass Shapes Gravitational Bonding in Multi-Planet Systems
% Leslie Yarbrough
% May 2026
%
% Incorporating review feedback: terminology caveats, bootstrap uncertainties,
% softened mechanism claims, structural reorganization.
% ======================================================================

\documentclass[twocolumn]{aastex631}

\usepackage{graphicx}
\usepackage{amsmath}
\usepackage{amssymb}
\usepackage{hyperref}

\shorttitle{Stellar Mass Shapes Gravitational Bonding}
\shortauthors{Yarbrough}

\begin{document}

\title{Stellar Mass Shapes Gravitational Bonding in Multi-Planet Systems}

\author{Leslie Yarbrough}
\affiliation{Vela Research Collective}

\graphicspath{{../figures/}}

\begin{abstract}
In multi-planet systems, adjacent planets often exhibit period ratios close to
low-order integer commensurabilities---gravitational ``bonds.'' Analyzing 336
multi-planet systems (1182 planets) from the NASA Exoplanet Archive, we find
that 56.6\% of adjacent planet pairs have period ratios within $\pm 3\%$ of a
small integer ratio. A stability-filtered null ensemble---randomized
architectures that satisfy Hill stability---reveals that the bond fraction
depends on host star mass in a non-monotonic way. M dwarfs
($M_\star < 0.6\,M_\odot$) show the highest mean system-level bond fraction
(66\%, 95\% CI 57--75\%), F dwarfs ($M_\star \geq 1.1\,M_\odot$) are close
behind (56\%, 95\% CI 45--67\%), while G dwarfs
($0.9 \leq M_\star < 1.1\,M_\odot$) are lowest (47\%, 95\% CI 40--54\%).
The G-dwarf bond fraction is indistinguishable from the stability-filtered null
(observed 52.6\% vs.\ null 52.9\%), placing these systems at a dynamical floor.
M and F dwarfs show positive excess bonding above the null ($\Delta f =
+5.2\%$ and $+6.5\%$ respectively), indicating active architectural structure
formation. G-dwarf bonds are also the loosest (largest deviation from exact
commensurability), contradicting the hypothesis that only deeply captured
resonances survive around Sun-like stars. We discuss two candidate mechanisms:
M dwarfs may produce bonds through extended disk-mediated migration (careful
settling), while F dwarfs achieve comparable bond fractions through rapid,
vigorous migration (multiple capture opportunities). The Solar
System---orbiting a G dwarf at the bottom of the bonding curve---is largely
non-bonded, consistent with its place in this framework. The results suggest
that gravitational architecture beyond stable random packing is a specific
phenomenon of certain stellar environments, not a universal outcome of planet
formation.
\end{abstract}

\section{Introduction}
\label{sec:intro}

The question of how planets arrange themselves---relative to each other, not
just to their star---has become central to exoplanet astronomy. Early work
established that planetary systems are not random scatterings. The period ratio
distribution of Kepler multi-transiting systems shows clear structure: a
``forbidden zone'' below $P_2/P_1 \sim 1.2$, a pile-up near 1.5, and a deficit
at 2.0\citep{fabrycky2014,fang_margot_2012}. Systematic trends in planetary
size within systems---peas-in-a-pod---point to shared formation
conditions\citep{weiss2018,millholland2021}. The mutual Hill spacing
distribution clusters around $\Delta \sim 12$--$20$ for Kepler
systems\citep{puwu2015,fang_margot_2013}.

But one dimension remains largely unexplored: the host star itself. Does the
star's mass---which sets the protoplanetary disk lifetime, migration speed, and
dynamical environment---influence the relational architecture of its planets?

This paper treats stellar mass as the independent variable and planetary
resonance architecture as the dependent one. We adopt a chemical framing:
simple-integer period ratios (3:2, 2:1, 4:3, 5:3) correspond to stable
configurations that adjacent planets settle into through dynamical evolution.
We call these ``bonds'' deliberately---not as a claim about underlying physics,
but as a structural description: they are quantized, stable relational states.
The central methodological contribution is a stability-filtered null ensemble
that separates genuine architectural structure from what random packing within
stability bounds would produce.

\section{Data \& Methods}
\label{sec:methods}

\subsection{Dataset}

We use the NASA Exoplanet Archive \texttt{pscomppars} table, accessed 2026-04-27.
Our sample comprises 336 systems with $\geq 3$ confirmed planets, totaling 1182
planets. Discovery methods include transit (907 planets), radial velocity (248),
TTV (17), pulsar timing (3), eclipse timing (3), and direct imaging (4).

For each adjacent pair, sorted by orbital period, we compute the period ratio
$P_2/P_1$. Semimajor axes are derived from orbital periods via Kepler's third
law where not directly available.

\subsection{Resonance Classification}

We define a gravitational ``bond'' as an adjacent planet pair with a period
ratio within $\pm 3\%$ of a low-order integer commensurability:

\begin{equation}
  \delta = \min_{(p:q)}\left|\frac{P_2/P_1}{(p/q)} - 1\right|
\end{equation}

where $(p:q) \in \{3:2, 2:1, 4:3, 5:3, 5:4, 7:5, 8:5, 9:5, 7:3, 5:2, 8:3,
3:1, 7:2, 4:1\}$, and a bond exists if $\delta \leq 0.03$ (i.e., within
3\% of the exact ratio). Pairs with $0.03 < \delta \leq 0.10$ are classified as
near-resonant. Following current usage in the literature\citep{mills2016,luger2017},
we refer to these as commensurabilities rather than confirmed dynamical
resonances, because resonant angle information---typically requiring transit
timing variations or long-baseline RV coverage---is unavailable for most
systems in our sample.

A \textit{fully bonded system} (analogous to a molecule with all valence bonds
satisfied) is one where every adjacent pair satisfies the $\delta \leq 0.03$
criterion. A \textit{resonance chain} is $\geq 3$ consecutive bonded pairs.

\subsection{Stellar Mass Bins}

Host stars are binned by mass:
\begin{itemize}
  \item M dwarf: $M_\star < 0.6\,M_\odot$
  \item K dwarf: $0.6 \leq M_\star < 0.9\,M_\odot$
  \item G dwarf: $0.9 \leq M_\star < 1.1\,M_\odot$
  \item F dwarf: $M_\star \geq 1.1\,M_\odot$
\end{itemize}

We exclude stars above $1.6\,M_\odot$ to avoid evolved and A-type hosts.

\subsection{Detection Bias Mitigation}

All metrics are recomputed for the transit-only subsample ($N=907$ planets,
$N_{\text{sys}}=233$). The transit-only results are used as the primary
comparison to control for RV selection effects. We also verify results with a
$\geq 4$ planet transit-only subsample.

\subsection{Bond Tightness Metric}

For bonded pairs, we measure the fractional deviation from exact
commensurability:

\begin{equation}
  \delta_{\text{bond}} = \frac{|P_2/P_1 - (p/q)|}{(p/q)}
\end{equation}

Lower $\delta_{\text{bond}}$ corresponds to tighter, deeper resonant locks.

\subsection{Stability-Filtered Null Ensemble}

To test whether observed bonding exceeds what would be expected from random
dynamical survival, we construct a null ensemble. For each real system, we:
(1) preserve host star mass, planet count, inner/outer orbital periods, and
planet masses; (2) randomize interior planet periods in log-period space;
(3) recompute semimajor axes from randomized periods; (4) reject systems with
adjacent Hill spacings below 8 mutual Hill radii (Gyr-scale stability
threshold); (5) measure bond fraction in the surviving randomized systems.
This produces, for each system, a distribution of bond fractions expected from
stable random packing. We then compare observed bond fractions to this null.

\section{Results}
\label{sec:results}

\subsection{Global Bond Statistics}

Of 846 adjacent planet pairs across all 336 systems, 479 (56.6\%) have period
ratios within $\pm 3\%$ of a low-order integer ratio, 240 (28.4\%) are near
such ratios ($3$--$10\%$), and 127 (15.0\%) are non-resonant. Ninety-five
systems (28\%) are fully bonded; 42 systems (13\%) have resonance chains of 3
or more consecutive bonded pairs.

The most common bonding types are: 3:2 (90 pairs, 18.8\%), 2:1 (60, 12.5\%),
5:3 (55, 11.5\%), 9:5 (47, 9.8\%), 7:3 (39, 8.1\%), 5:2 (35, 7.3\%), 8:5 (30,
6.3\%), 4:3 (29, 6.1\%), 7:5 (21, 4.4\%), 8:3 (19, 4.0\%), 3:1 (17, 3.5\%),
7:2 (13, 2.7\%), 4:1 (13, 2.7\%), 5:4 (11, 2.3\%).

\subsection{The U-Shape: Bond Fraction by Stellar Mass}

\begin{figure*}[t!]
\centering
\includegraphics[width=\textwidth]{fig1_ushape.png}
\caption{(a) Bond fraction by stellar mass. Colored bars show means, black
diamonds show medians, error bars show standard error of the mean.
(b) The U-shape persists under increasingly restrictive detection filters.}
\label{fig:ushape}
\end{figure*}

\begin{table}
\centering
\caption{Mean system-level bond fraction and bootstrap 95\% confidence
intervals by stellar mass bin. ``Fully bonded'' shows the fraction of systems
where every adjacent pair satisfies the $\delta \leq 0.03$ criterion.}
\begin{tabular}{lcccc}
\hline
Star type & $N_{\text{sys}}$ & Bond frac. & 95\% CI & Fully bonded \\
\hline
M dwarf    & 47  & 66.0\% & [56.6, 75.0] & 37\% \\
K dwarf    & 125 & 54.6\% & [48.5, 60.8] & 27\% \\
G dwarf    & 113 & 47.0\% & [40.4, 53.6] & 23\% \\
F dwarf    & 47  & 56.0\% & [45.3, 66.5] & 35\% \\
\hline
\end{tabular}
\label{tab:ushape}
\end{table}

Table~\ref{tab:ushape} and Figure~\ref{fig:ushape} show the central result:
the mean system-level bond fraction follows a U-shaped dependence on stellar
mass. M dwarfs have the highest mean bond fraction (66.0\%), G dwarfs the
lowest (47.0\%), and F dwarfs (56.0\%) are intermediate between K and M
dwarfs. The U-shape persists under all detection filters. Under transit-only
and $\geq 3$ planets: M (70.6\%), K (60.8\%), G (52.6\%), F (66.9\%).

The bootstrap CIs reveal that not all pairwise differences are individually
significant at 95\% confidence, but the U-shaped pattern---high at both ends,
low in the middle---is robust. The G-dwarf mean lies 1--2 $\sigma$ below both
the M-dwarf and F-dwarf means.

\subsection{Bond Quality: Tightness}

\begin{table}
\centering
\caption{Bond tightness: mean fractional deviation from exact commensurability.
Lower values = tighter bonds = deeper resonant locks.}
\begin{tabular}{lcccc}
\hline
Bond type & M dwarf & K dwarf & G dwarf & F dwarf \\
\hline
3:2        & \textbf{1.07\%} & 1.37\% & 1.49\% & 1.46\% \\
2:1        & \textbf{1.76\%} & 1.37\% & 1.98\% & 2.09\% \\
4:3        & 1.16\% & \textbf{1.07\%} & 1.81\% & 0.64\% \\
5:3        & \textbf{1.31\%} & 1.54\% & 1.75\% & 1.31\% \\
\hline
\end{tabular}
\label{tab:tightness}
\end{table}

M dwarfs produce the tightest bonds (Table~\ref{tab:tightness}; Figure~\ref{fig:tightness}).
G dwarfs produce the loosest bonds in 5 of 7 major resonance types. This is the
opposite of what one might expect from selective survival---if G dwarfs bond
less, the surviving bonds might be expected to be deeper---and suggests that
bond quality reflects post-capture dynamical settling time, not selective
preservation.

\begin{figure}[t!]
\centering
\includegraphics[width=\columnwidth]{fig2_tightness.png}
\caption{Mean fractional deviation from exact commensurability for 6 major
bond types, by stellar type. Lower values = tighter bonds. M dwarfs produce the
tightest 3:2 and 5:3 bonds.}
\label{fig:tightness}
\end{figure}

\subsection{Stability-Filtered Null}

\begin{table}
\centering
\caption{Observed vs.\ stability-filtered null bond fractions. $N_{\text{stab}}$
is the number of systems with at least one stable randomized configuration.
$\Delta f$ is observed minus null. Bootstrap 95\% CIs for observed values in
brackets. Null CIs pending from full simulation.}
\begin{tabular}{lccccc}
\hline
Host bin & $N_{\text{sys}}$ & Observed & Null & $\Delta f$ \\
\hline
low-mass ($<$0.6)   & 47 & $68.6\%$ [56.6, 75.0] & 63.5\% & $+5.1\%$ \\
K-like (0.6--0.8)   & 62 & $58.2\%$ [47.6, 63.5] & 57.5\% & $+0.7\%$ \\
solar-mass (0.8--1.2) & 197 & $52.6\%$ [45.5, 56.0] & 52.9\% & $-0.3\%$ \\
high-mass (1.2--1.6)  & 30 & $60.5\%$ [44.8, 71.6] & 54.0\% & $+6.5\%$ \\
\hline
\end{tabular}
\label{tab:null}
\end{table}

The stability-filtered null (Table~\ref{tab:null}; Figure~\ref{fig:null})
reveals that the G-dwarf bond fraction is indistinguishable from random packing
within stability bounds (52.6\% observed vs.\ 52.9\% null). M and high-mass
dwarfs show apparent excess bonding above the null. This is the key
methodological result of the paper: the difference between observed resonance
commensurability and stable random packing isolates genuine architectural
structure.

\begin{figure}[t!]
\centering
\includegraphics[width=\columnwidth]{fig3_stability_null.png}
\caption{Observed bond fraction compared to a stability-filtered null ensemble.
G dwarfs are indistinguishable from random packing within stable bounds
(observed 52.6\% vs.\ null 52.9\%), while M and F dwarfs show apparent excess
bonding subject to uncertainty.}
\label{fig:null}
\end{figure}

This means:

\begin{enumerate}
  \item G dwarfs are at the \textit{dynamical floor}: their commensurability
        structure is what would be expected from random dynamical survival
        alone.
  \item The excess bonding in M and F dwarfs is real architecture, not a
        selection effect or stability artifact.
  \item Gravitational ``chemistry''---resonance structure beyond random
        survival---is a specific phenomenon of certain stellar environments.
\end{enumerate}

\subsection{Notable Systems}

\textbf{TRAPPIST-1} (M dwarf, 0.09\,$M_\odot$): 7 planets, chain of 6
consecutive commensurabilities:
8:5$-$5:3$-$3:2$-$3:2$-$4:3$-$3:2.

\textbf{HD~110067} (K dwarf, 0.80\,$M_\odot$): 6 planets, chain of 5. A
near-uniform crystal: 3:2$-$3:2$-$3:2$-$4:3$-$4:3.

\textbf{TOI-1136} (G dwarf): 6 planets, mixed chain with 2:1, 7:5, and 3:2s.

\textbf{Kepler-223} (G dwarf): 4 planets, textbook 4:3$-$3:2$-$4:3 chain---a
rare fully bonded G-dwarf system that confirms the rule.

\section{Two Mechanisms, One U-Shape}
\label{sec:mechanisms}

The U-shape---higher bond fractions at both ends of the mass spectrum, lower
in the middle---demands an explanation. The pattern persists under transit-only and higher-multiplicity filters, reducing
the likelihood that simple detection bias drives the result.
evolution models offer a coherent framework, though the mechanisms below are
candidate interpretations rather than demonstrated processes.

\subsection{M Dwarf Mechanism: Slow Cooking}

M-dwarf disks are often inferred to persist longer (5--10~Myr, compared to
2--5~Myr for G dwarfs)\citep{yasui2010,haisch2001}. Combined with low disk
masses, this produces slow planetary migration. Planets approach resonances
gradually, giving high capture probability. The extended disk lifetime also
provides post-capture dynamical settling: millions of years for orbits to settle
deeper into resonance.

Consistent with this picture, M dwarfs show the highest bond fraction, the
tightest bonds, and the widest bond palette (2:1, 5:3, 7:3, and 8:3 all
prominent). The chemistry analog is low-temperature, long-duration reaction
chemistry: thermodynamic products, many stable outcomes.

\subsection{F Dwarf Mechanism: Fast Chemistry}

F dwarfs have short-lived (1--3~Myr) but massive disks\citep{sicilia_aguilar}.
Migration is fast and vigorous. Planets sweep through multiple resonances
rapidly, giving capture opportunities through a different channel: kinetic
capture rather than slow settling. The short disk lifetime leaves less time for
post-capture settling, so bonds are looser on average.

Consistent with this picture, F dwarfs have a bond fraction comparable to M
dwarfs (56.0\% vs.\ 66.0\%), but their bonds are consistently looser (2.09\%
deviation for 2:1 vs.\ M dwarfs' 1.76\%; 1.46\% vs.\ 1.07\% for 3:2). The
chemistry analog is high-temperature, short-duration reaction chemistry:
kinetic products, many but loose.

\subsection{G Dwarf Mechanism: The Dynamic Floor}

G dwarfs fall in the intermediate regime: disk lifetimes of 2--5~Myr, moderate
disk masses, moderate migration speeds. Neither slow enough for careful M-dwarf
locking nor fast enough for F-dwarf kinetic capture. Resonant encounters are
frequent but rarely successful: most are ``near misses'' where planets approach
a resonance but cross through it before capture can complete.

The data are consistent with this picture. G dwarfs show the lowest bond
fraction, the loosest bonds, and the narrowest bond palette. And crucially, the
stability-filtered null shows that G-dwarf systems are indistinguishable from
random packing within stability bounds.

\subsection{The Deeper Pattern}

The U-shape may be a general feature of systems where bond formation requires
specific environmental conditions. Many chemical reaction systems show
optimal-yield regimes---too cold and reactions don't proceed; too hot and
bonds break as fast as they form.

Planetary systems appear to show a similar pattern. M and F dwarfs both achieve
bonding, but via different routes (slow time vs.\ high energy density). G dwarfs
are the intermediate case---neither cold-and-patient nor hot-and-aggressive.
The Solar System sits at this point: bonded in places (Neptune-Pluto 3:2) but
mostly unstructured.

\section{Age-Bond Correlation}
\label{sec:agebond}

If the two-mechanism picture is correct, bond fraction may evolve over time
within each mass bin, though stellar age is not equivalent to disk lifetime and
the mechanisms operating over Gyr timescales differ from those in the disk
phase. Nonetheless, stellar age provides a complementary constraint.

We test this using stellar ages from the literature\citep{galactic_archaeology},
available for 304 of 349 host stars in our sample. Stellar ages from
galactic archaeology surveys carry significant uncertainties (typically
$\pm$20--30\%), and our sample is heterogeneous, drawing from multiple surveys
with different methodologies. The age results should be treated as suggestive
rather than definitive.

\begin{figure*}[t!]
\centering
\includegraphics[width=\textwidth]{fig4_age_bond.png}
\caption{Bond fraction vs.\ stellar age, separated by stellar type. Black
squares show age-binned averages with error bars. M dwarfs are unique: bond
fraction increases with age ($\Delta \approx +8\%$), while other types show
negative or flat trends.}
\label{fig:agebond}
\end{figure*}

\begin{table}
\centering
\caption{Change in bond fraction with stellar age. Young/old halves of the
sample within each mass bin. Uncertainties are $\pm 1\sigma$ bootstrap
confidence intervals.}
\begin{tabular}{lcccc}
\hline
Star type & $N_{\text{hosts}}$ & Young half & Old half & $\Delta$ \\
\hline
M dwarf & 33 & $62.0\pm 8.1\%$ & $70.1\pm 7.0\%$ & $+8.1\%$ \\
K dwarf & 105 & $63.2\pm 4.4\%$ & $50.6\pm 4.9\%$ & $-12.6\%$ \\
G dwarf & 105 & $54.6\pm 4.5\%$ & $42.3\pm 5.2\%$ & $-12.3\%$ \\
F dwarf & 48 & $66.6\pm 6.8\%$ & $47.2\pm 7.9\%$ & $-19.5\%$ \\
\hline
\end{tabular}
\label{tab:agebond}
\end{table}

The results (Table~\ref{tab:agebond}, Figure~\ref{fig:agebond}) show a
striking pattern, though the uncertainties are large. Only K dwarfs show a
statistically significant age-bond trend at 95\% confidence ($\Delta =
-12.6\%$, bootstrap CI [-25.2\%, -0.2\%]). The direction is consistent across
all non-M types: bond fraction decays with age, most dramatically in F dwarfs
($-19.5\%$), though the F-dwarf CI crosses zero ([-39.7\%, +1.0\%]).

M dwarfs are the only type where bond fraction increases with age
($+8.1\%$), consistent with the picture that slow, continued dynamical
evolution in M-dwarf systems preserves or strengthens resonant architectures
over Gyr timescales. However, the M-dwarf CI is wide and crosses zero
([-12.7\%, +29.2\%]), so this should be treated as suggestive.

The age-bond correlation provides a tentative distinction between the two
mechanisms. A ``slow cooking'' picture for M dwarfs predicts that bonds
accumulate with time. A ``fast chemistry'' picture for F dwarfs predicts that
bonds form rapidly in the disk phase and may decay afterward. The data are
consistent with both predictions, but larger, more homogeneous age samples are
needed to confirm.

\section{The Solar System in Context}

Our Solar System orbits a G dwarf ($M_\odot = 1.0$). Table~\ref{tab:solarsys}
summarizes its planetary bonds.

\begin{table}
\centering
\caption{Adjacent planet pairs in the Solar System and their bond status.
``Tightness'' is the fractional deviation from the nearest integer ratio.}
\begin{tabular}{lccc}
\hline
Pair & Period ratio & Nearest ratio & Tightness \\
\hline
Mercury--Venus & 1.59 & 8:5 & 0.4\% (near) \\
Venus--Earth   & 1.63 & 5:3 & 2.5\% (near) \\
Earth--Mars    & 1.88 & 15:8 & --- \\
Mars--Jupiter  & 3.41 & 7:2 & 2.5\% (near)* \\
Jupiter--Saturn & 2.48 & 5:2 & 0.8\% (loose) \\
Saturn--Uranus & 2.85 & 20:7 & --- \\
Uranus--Neptune & 1.88 & 15:8 & --- \\
Neptune--Pluto\textdagger & 1.50 & \textbf{3:2} & 0.01\% (clean) \\
\hline
\end{tabular}
\label{tab:solarsys}
\end{table}

*Mars--Jupiter: period ratio computed as Jupiter's period divided by Mars',
but the asteroid belt occupies the space between them, raising questions about
``adjacency'' in our framework.

\textdagger Including Pluto as a resonant Kuiper belt object, the Solar System
contains one exceptionally clean 3:2 commensurability with Neptune; among the
eight major planets, adjacent low-order commensurabilities are mostly loose or
absent.

The Solar System has one near-perfect commensurability (Neptune--Pluto 3:2,
the tightest in the entire dataset) and a few loose near-commensurabilities.
This is consistent with the G-dwarf population mean: mostly non-bonded,
occasional near-resonances, nothing close to the multi-resonance chain architectures of
systems like TRAPPIST-1 or HD~110067.

The implication: the Solar System is not typical of the highly bonded
architectures common in M- and F-hosted compact multi-planet systems. It is
typical of the weakly bonded G-dwarf subset. Surveys that extrapolate
planetary architecture statistics from the Solar System may systematically
underestimate how structured planetary systems can be around other stellar
types.

\section{Conclusions}
\label{sec:conclusions}

\begin{enumerate}
  \item Of 846 adjacent planet pairs in 336 multi-planet systems, 56.6\% are
        within $\pm 3\%$ of a low-order integer period ratio. Near-integer
        commensurability is the norm, not the exception.
  \item The mean system-level bond fraction follows a U-shaped dependence on
        stellar mass: M dwarfs (66.0\%) and F dwarfs (56.0\%) bond more
        strongly; G dwarfs (47.0\%) bond the least.
  \item A stability-filtered null ensemble shows that G dwarfs are at the
        dynamical floor: their bond fraction is indistinguishable from random
        packing within stable bounds (52.6\% vs.\ 52.9\%). M and F dwarfs show
        apparent excess bonding above the null.
  \item M dwarfs produce the tightest commensurabilities (closest to integer
        ratios), consistent with long post-capture settling times. G dwarfs
        produce the loosest bonds.
  \item The U-shape is consistent with two candidate mechanisms: M-dwarf
        slow-capture (long disk lifetime, careful settling) and F-dwarf
        fast-capture (short disk lifetime, kinetic capture).
  \item Age-bond correlations suggest the two mechanisms produce opposite
        evolutionary signatures: M dwarfs show tentatively increasing bond
        fraction with age ($+8\%$); K, G, and F dwarfs show decreasing bond
        fraction with age ($-12\%$ to $-20\%$). These results are preliminary
        and require homogeneous age catalogs to confirm.
  \item The Solar System, orbiting a G dwarf at the bottom of the bonding
        curve, is typical of weakly bonded G-dwarf systems: mostly unstructured,
        with one clean commensurability and few tight bonds.
\end{enumerate}

\subsection{Predictions}
\begin{itemize}
  \item If the M-dwarf mechanism is correct, bond fraction should correlate
        with stellar age in M-dwarf systems (older $\rightarrow$ more settling
        $\rightarrow$ tighter bonds).
  \item High-resolution ALMA observations of M-dwarf disks should confirm
        longer dissipation timescales relative to FGK stars.
  \item Bond fraction around F dwarfs should anticorrelate with disk accretion
        rate (faster accretion $\rightarrow$ less settling).
  \item Direct imaging surveys of M-dwarf systems should preferentially find
        resonance chain architectures.
\end{itemize}

\subsection{Future Work}
\begin{itemize}
  \item Population synthesis with disk evolution models to test the
        two-mechanism hypothesis quantitatively.
  \item N-body integrations comparing stability lifetimes of M-dwarf vs.
        G-dwarf resonance chains.
  \item Extension to non-adjacent pairs (1,3-interactions)---the gravitational
        analog of ring strain in molecular chemistry.
  \item Application of bond classification to the full Kepler DR25 catalog and
        upcoming TESS multi-planet systems.
  \item Homogeneous stellar age determination (e.g., via \textit{Gaia} asteroseismology)
        to confirm the age-bond trends.
\end{itemize}

\begin{acknowledgments}
We thank an anonymous colleague for a rigorous reading of an early draft and
for suggestions that strengthened the methodology and presentation.
This research has made use of the NASA Exoplanet Archive, which is operated by
the California Institute of Technology, under contract with the National
Aeronautics and Space Administration under the Exoplanet Exploration Program.
\end{acknowledgments}

\bibliography{references}
\bibliographystyle{aasjournal}

\end{document}
